\section{Comparison with the tiling-specific automata}
\label{sec:comparison}

It is worth setting this construction beside Margenstern's, since the two answer different questions.

\subsection{A family of automata against one}

The tiling-specific automata are a family, a separate object per tiling, each with its own table, summarized in
Table~\ref{tab:compare}. Adding a tiling means designing a new automaton and checking a new table. The present
construction is one automaton. Adding a tiling means nothing, since the rules already apply to its cell graph.
The family covers the tilings it has reached, the single automaton covers every regular tiling, including the
ones never individually reached, such as the octagrid.

\begin{table}[ht]
\centering
\begin{tabular}{@{}lllc@{}}
\toprule
construction & tiling & states & rotation invariant \\
\midrule
Margenstern, 5-state pentagrid \cite{margenstern2014penta5} & $\{5,4\}$ & 5 & yes \\
Margenstern, 3-state pentagrid \cite{margenstern2015penta3} & $\{5,4\}$ & 3 & yes \\
Margenstern, 2-state pentagrid \cite{margenstern2015penta2} & $\{5,4\}$ & 2 & yes \\
Margenstern, heptagrid \cite{margenstern2009hepta} & $\{7,3\}$ & 4 & yes \\
Margenstern, dodecagrid \cite{margenstern2021dodeca5} & $\{5,3,4\}$ & 5 & yes \\
Margenstern, totalistic dodecagrid \cite{margenstern2021dodeca4} & $\{5,3,4\}$ & 4 & yes \\
\midrule
this paper (weak) & every regular tiling & 4 dynamic + switch bit & no \\
this paper (strong) & every regular tiling & 4 dynamic + switch bit + build & no \\
\bottomrule
\end{tabular}
\caption{The tiling-specific automata against the uniform one. Each of Margenstern's is rotation invariant and
small but solves a single tiling. Ours is not rotation invariant, since its switches are oriented, but one
construction covers the whole family, and the strong version starts from a finite seed.}
\label{tab:compare}
\end{table}

\subsection{Rotation invariance against oriented switches}

The deeper difference is how each recovers direction without a global frame. Margenstern's automata are rotation
invariant, a rule holds under every rotation of the cell, and direction is carried in the states with the help
of milestones. Rotation invariance is strong and elegant and is the reason those automata read as cellular
automata in the strict sense, but the way a rotation permutes a cell's neighbours differs in each tiling, so it
must be re-established for every tiling, which is much of why each needs its own table.

The present automaton is not rotation invariant. Its switches are oriented, carrying labels for the trunk and
the two branches in the initial configuration. Because the labels are local data, not a global frame, the
automaton stays tiling-independent, but it is no longer rotation invariant. The track network is the program,
and the labels are part of the program.

\subsection{The locomotive's direction trick}

One point of technique lets the dynamic alphabet stay small while remaining tiling-independent. On a bare graph
nothing tells the head which way is forward. The two-cell locomotive supplies the answer at no cost, by
Lemma~\ref{lem:direction} the cell behind the head is the trailing cell and the only clear track cell adjacent
to the head is the one ahead. Forward is read off the local pattern, not the plane. This is the service
Margenstern's milestones provide, achieved here by the shape of the locomotive, which is why no decoration of
the track and no orientation of plain track cells is needed.
